\section{Simulation Evaluation}
\label{sec_eval}
\subsection{Questions and Method}
We ask three questions: (Q1) does the simulated two-signal controller reduce
footprint variability relative to no controller; (Q2) how sensitive are outcomes
to $\alpha$, $\theta_1$, and $\theta_2$; and (Q3) can every reported value be
regenerated from a seed and script?  We do not ask whether HMS outperforms Android
on a phone.

The artifact generates four synthetic traces: W1 (bursty camera-like allocation),
W2 (streaming game assets), W3 (inference-like tensor bursts), and W4
(multi-application switching).  A run is identified by configuration and random
seed.  The baseline and HMS consume the same generated trace.  Commands, output
schema, and provenance are documented in the repository README.  Values are
simulation units scaled for readable plots; they are not milliseconds, megabytes,
watts, or device confidence intervals unless independently calibrated.

\subsection{Metrics}
Heap stability has one definition throughout the paper and code:
\[
S_h=1-\frac{\sigma(R)}{\max(\mu(R),\epsilon)} .
\]
The evaluation window is the complete post-warm-up trace.  We report $S_h$ together
with peak RSS because a stable but large footprint should not be rewarded without a
size metric.  $L_r$ is simulated reclaim-service latency, never ``memory seek
latency.''  $T_{gc}$ is simulated pause duration.  Tail percentiles, deadline
misses, energy, and effect sizes require per-run device data and are not claimed.

\begin{figure*}[t]
\centering
\begin{subfigure}{0.42\textwidth}\centering
  \includegraphics[width=\linewidth]{fig9_W1_alpha_vs_Sh_3alpha.png}
  \caption{W1: $\alpha$ versus $S_h$}
  \Description{A line plot of simulated W1 stability over an alpha sweep.}
\end{subfigure}\hspace{0.04\textwidth}
\begin{subfigure}{0.42\textwidth}\centering
  \includegraphics[width=\linewidth]{fig9_W2_theta1_vs_Lr.png}
  \caption{W2: $\theta_1$ versus $L_r$}
  \Description{A line plot of simulated W2 reclaim latency over a threshold sweep.}
\end{subfigure}
\begin{subfigure}{0.42\textwidth}\centering
  \includegraphics[width=\linewidth]{fig9_W3_theta2_vs_Sh.png}
  \caption{W3: $\theta_2$ versus $S_h$}
  \Description{A line plot of simulated W3 stability over a threshold sweep.}
\end{subfigure}\hspace{0.04\textwidth}
\begin{subfigure}{0.42\textwidth}\centering
  \includegraphics[width=\linewidth]{fig9_W4_theta2_vs_Lr.png}
  \caption{W4: $\theta_2$ versus $L_r$}
  \Description{A line plot of simulated W4 reclaim latency over a threshold sweep.}
\end{subfigure}
\caption{Seeded simulation parameter sweeps.}
\Description{Four plots show sensitivity of simulated stability and reclaim
latency to controller weights and thresholds.}
\label{fig:parameter_sensitivity}
\end{figure*}

\subsection{Interpretation}
The simulation exposes a mixed trade-off rather than a universal improvement:
with the documented seed, controller enforcement reduces the reclaim- and
GC-service proxies and peak RSS, while $S_h$ can decrease because the bounded
interventions add footprint variation.  Parameter sweeps also show that the outcome
depends on threshold selection.  This is a sanity-check of controller logic and a
falsifiable warning against relying on a single metric, not empirical validation.
Exact numerical
claims from earlier drafts---including commercial-device p99 latency,
misprediction rate, post-reclaim faults, compatibility, CPU overhead, and
low-pressure regression---have been removed because their raw inputs are absent.

\subsection{Required Device Evaluation}
A future empirical paper must publish one row per device, workload, system, and run
with peak RSS; GC p50/p95/p99; reclaim p99; frame p99 and miss ratio; energy;
reclaim volume; refaults; and application kills.  Raw Perfetto traces,
\texttt{hms\_events.jsonl}, RSS, vmstat, PSI, configuration, exclusions, and analysis
scripts must accompany those rows.

Required baselines are stock Android, tuned LMKD/PSI, direct
\texttt{memory.reclaim}, equal-volume static reclaim, a PSI-driven userspace
daemon, DAMON/DAMOS, MGLRU tuning, footprint-only, rate-only, and signal ablations.
At least four memory tiers and two SoC vendors should be studied.  Long-duration,
energy, thermal, writeback, functional-correctness, and multi-process fairness
tests are also required.

Statistical analysis must be selected after inspecting the raw distribution.
Repeated measures should preserve pairing; device and run effects should not be
pooled as independent observations.  Bootstrap intervals or paired nonparametric
tests are appropriate for heavy-tailed metrics, with Holm correction and an effect
size such as Cliff's delta.  No t-test or confidence interval is reported here
because the repository contains no qualifying device observations.
